\section{Full transfer ordered by quantum grading}
\label{sec:graded-transfer}

The graded kernel from report 23 supplies another exact response to slow
sparse cancellation.  It computes all maps between scalar survivors,
including adjacent-degree maps, without the common-disk certificate required
by Section~\ref{sec:disk-transfer-integration}.  Production exposes \code{graded} and
\code{graded-adaptive} as optional reduction policies.  This integration
covers the graded kernel; the report's later twist-threshold, succinct CLI,
and geometric-dependency proposals have separate review obligations.

\subsection{Retaining the grading rather than reconstructing it afterwards}
At a frontier of $2k$ points let $c(a,b)$ be the number of circles in the
overlay of matchings $a,b$.  A basis morphism dotted on a set $S$ of those
circles has intrinsic weight
\[
 \omega(a,b,S)=k-c(a,b)+2|S|.
\]
The existing coefficient engine and grading audit use the convention
$q_b-q_a=\omega$ for every nonzero monomial in a homogeneous differential
entry.  The new scanner keeps these object shifts during allocation:
\[
 q'=q+i+c-2\operatorname{popcount}(\ell),
\]
where $i\in\{0,1\}$ is the smoothing, $c$ is the number of delooped circles,
and $\ell$ is their binary label.  The array follows the exact object order
used by \code{FastScan.add\_crossing}.  Scalar cancellation preserves the
shifts of all remaining objects.  This is the ordinary dotted grading and
cancellation calculus of \cite{barnatan}, used here as an algorithmic order.

Weights are nonnegative because $c(a,b)\le k$.  Equality requires equal
matchings and no dots: a circle in the overlay then consists of the two
copies of the same arc.  Composition adds weights on every nonzero term.
In the engine's dot-count convention, for a composition through $b$ with
input dot counts $u,v$ and output count $r$,
\[
 2r=2u+2v+k-c(a,b)-c(b,c)+c(a,c),
\]
which proves additivity by substitution.  Thus writing $d=d_0+\delta$,
the zero-weight part $d_0$ is a direct sum of binary complexes at fixed
matching and $q$, and $\delta$ strictly raises $q$.  An undotted morphism
between distinct matchings is not a scalar identity.  Taking weight zero
in $d^2=0$ proves $d_0^2=0$.

\subsection{A contraction with representatives and homotopies}
For each fixed matching and quantum shift, binary elimination splits each
chain group as
\[
 C^h=B^h\oplus R^h\oplus L^h,
\]
where $B^h=\operatorname{im}d_0^{h-1}$, $R^h$ represents its quotient
in $\ker d_0^h$, and $d_0:L^h\to B^{h+1}$ is an isomorphism.
Let $I,P$ include/project $R$ and let $H$ invert the last isomorphism,
vanishing on the other summands.  Then
\[
 PI=1,\quad d_0I=Pd_0=0,\quad d_0H+Hd_0=1+IP,
 \quad H^2=HI=PH=0.
\]
The implementation constructs bases, lifts, and inverse coordinates; a
list of ranks would not supply the maps needed for this step.  Each scalar
map preserves matching and $q$, and $H$ lowers homological degree by one.

There are finitely many object shifts, so $H\delta$ and $\delta H$ are
nilpotent.  In characteristic two form the finite inverses
\[
 U=(1+H\delta)^{-1},\qquad V=(1+\delta H)^{-1}.
\]
The finite perturbation identity \cite{crainic-perturbation} gives
\begin{equation}
 D=P\delta UI,\quad I'=UI,\quad P'=PV,\quad H'=UH=HV.
 \label{eq:graded-full-transfer}
\end{equation}
These are a special chain contraction from $(C,d)$ to $(R,D)$:
\[
 dI'=I'D,\quad P'd=DP',\quad P'I'=1,\quad
 dH'+H'd=1+I'P',\quad H'^2=H'I'=P'H'=0,\quad D^2=0.
\]
For an explicit algebraic check, $Z=UI$ solves $Z=I+H\delta Z$.
Using $d_0\delta+\delta d_0+\delta^2=0$ and the scalar contraction gives
$dZ=ID+H\delta dZ$, hence $dZ=ZD$.  The right-hand calculation gives
the projection identity.  Expanding the finite inverses with $H^2=HI=PH=0$
proves the side identities.  Multiplying
$dH'+H'd+1+I'P'$ on the left by $1+H\delta$ and on the right by
$1+\delta H$ reduces it to zero using $d^2=0$ and the scalar homotopy
equation.  The multipliers are invertible.  Finally
$D^2=P'd^2I'=0$.  Every term of $D$ contains $\delta$ and has positive
weight, so the returned complex has no scalar unit left to cancel.

This termination argument uses the actual finite grading and typed
composition law.  It does not invoke a uniform nilpotence claim for
arbitrary ungraded formal matching tables.  Both the disk-certified and
graded algorithms remain applicable only to valid complexes in their
respective contracts.

\subsection{Computing a path sum without enumerating its paths}
For each scalar survivor column $e$, evaluate
\[
 Z=Ie+HW,\qquad W=\delta Z,\qquad De=PW
\]
one quantum block at a time in increasing order.  All contributions to a
block arrive from lower blocks before it is processed; binary $P,H$ then
mix within that block.  Each radical edge is propagated at most once for
this source, combining equal destinations by XOR.  This computes the
finite sum in \eqref{eq:graded-full-transfer} without constructing matrix
powers or enumerating exponentially many individual paths.

There is also a source-dependent band:
\[
 0\le\omega=k-c(a,b)+2|S|\le k+c(a,b)\le2k.
\]
A source of shift $q_0$ cannot reach a nonzero coefficient beyond $q_0+2k$.
The kernel drops such destinations before composition.  It still retains
higher-order corrections within the band.  For example, two radical steps
with a scalar contracting arrow between them can yield $x_1x_2$ even when
$P\delta I=0$.  On a single arc the analogous $x^2$ term is zero.  Both
cases are explicit regression tests, as is the nonzero minimal map
$a\xrightarrow{x}a$.  Replacing the latter by two isolated survivors would
change later gluing.

Let $N$ count allocated object slots including holes left by completed
sparse pivots, $S$ count scalar survivors, and $E_\delta$ count positive-weight
entries.  There are at most $S E_\delta$ coefficient-composition attempts.
If $A_H,A_P$ count nonzero scalar incidences, the remaining transfer work is
bounded by $O(S(N+E_\delta+A_H+A_P))$ coefficient additions and control
operations, after polynomial binary preprocessing.  Packed indices and
integer costs must also be charged.  The implementation's global bit columns
and arrays can depend on $N$, even if few slots remain live.  Thus a live-only
bound must first account for compaction or the size of the allocated arrays.
The component coefficient engine costs $2^{O(k)}\operatorname{poly}(k)$ per
operation, giving an explicit bound polynomial in $N,S$ and exponential in
$k$.  None of these statements bounds necessary object multiplicity.

Optional diagnostic mode additionally computes all of $I',P',H'$ and checks
the eight displayed identities in the full cobordism algebra.  It is excluded
from timed production reductions.  If the whole differential is scalar and
diagnostics are disabled, the kernel obtains quantum-resolved homology
multiplicities directly and avoids constructing unused contraction maps.

\subsection{Adaptive scheduling and installation}
The eager policy transfers each stage.  The adaptive policy first allows
$\max(256,4M)$ candidate Schur update pairs, where $M$ is the live count at
stage entry.  The existing sparse eliminator pauses only between complete
pivots and clears its candidate queues.  Full transfer then consumes that
partially reduced complex, retaining completed cancellations.  Counters
distinguish stages, actual switches, coefficient attempts, identity shortcuts,
and scalar-only reductions.

The maintained adapter prepares the new incoming-edge index before replacing
any differential arrays, checks the shared deadline again, and installs the
matching, degree, quantum, outgoing and incoming arrays together.  A failure
while preparing the replacement preserves the current differential; composition
caches may have acquired reusable entries.  Tests force both a final-boundary
deadline expiration and an incoming-index allocation failure, then resume the
same object successfully.  The kernel also polls during long composition
sequences.  A global resource failure propagates from the raw API and becomes
UNKNOWN in recognition; no rate of progress is a topological verdict.

Direct constructors have an optional transfer-object ceiling; exceeding it
resumes ordinary cancellation.  The public reduction modes retain the existing
object and time ceilings.  Their policy is a heuristic, not a competitive
running-time guarantee or a byte-memory budget.  In particular, a dense binary
basis change can be more expensive than the sparse pivots it replaces and can
change future fill.  The default remains \code{standard}.

Both modes require minimum-fill pivots, bit coefficients, self-inverse
cancellation, and one scan order.  Standard and component coefficient
composition and tail finishing are supported.  Homological windows are
rejected before early recognition certificates: their pruned allocation needs
a separate quantum-label adapter.  Shared, saturated, twist and other
alternative scanners are also rejected.  Existing disk-adaptive window modes
remain available.

\subsection{Validation and measurement scope}
The delivered report 23 snapshot passes all 281 tests unchanged and its
114 source Git blobs match the declared baseline.  Those full archive tests
also cover additions outside this integration; passing them does not constitute
integration of those additions.  Production adds 13 graded integration tests
and two transactional-resource tests to the previous 411-test suite.
The transfer audit uses 86 diagrams in 258 order presentations and checks
1812 full stage contractions, with eight identities each.  Four scanner
variants agree on raw ranks and three quantum-tracking variants agree on
quantum-resolved ranks.  Synthetic controls retain nonzero minimal maps,
test higher corrections, exercise capacity fallback and resume after sparse
progress.  Independent set-coefficient scans provide an additional oracle
on 18 small random diagrams.

The maintained drivers are in \code{fast/graded\_research}; raw audit and
source-provenance evidence is in \code{synthesis/data/graded-*}.
Measurements use seven shuffled paired rounds, an identical standard A/A arm,
fresh states and batched calls.  Whole-diagram timings cover complete raw scans
including reduction setup, but exclude construction and order selection;
they bypass recognition filters.  Synthetic timings cover the entire reduction
stage and exclude cloning.  Synthetic controls are valid graded algebraic
complexes, not proved realizations of classical knot prefixes.

\begin{center}\small
\begin{tabular}{lrrrr}
\toprule
Case & Old/eager & Old/adaptive & A/A & Adaptive switches\\
\midrule
hard\_unknot\_8 & 0.689 & 0.957 & 0.998 & 0\\
conway & 0.557 & 0.920 & 0.929 & 0\\
kinoshita\_terasaka & 0.616 & 0.980 & 1.019 & 0\\
unknot\_braid40 & 0.633 & 0.886 & 0.984 & 0\\
stress\_braid5\_36 & 0.260 & 0.942 & 1.021 & 0\\
synthetic 16 & 0.888 & 0.741 & 0.999 & 1\\
synthetic 32 & 1.909 & 1.582 & 0.992 & 1\\
synthetic 64 & 4.968 & 4.174 & 1.008 & 1\\
synthetic 128 & 13.180 & 11.266 & 1.015 & 1\\
\bottomrule
\end{tabular}
\end{center}
Ratios above one favor the new policy. Each is a median of paired ratios.
The five actual diagrams invoke zero adaptive transfers. Eager transfer slows
all five raw scans; tracking the grading also adds overhead to adaptive scans.
The largest synthetic case gains $13.18\times$ eagerly and $11.27\times$
adaptively, retaining a nonzero dotted minimal differential. The smallest
synthetic case regresses. A/A ratios range from 0.929 to 1.021 across all
cases, so small changes should not be treated as reliable improvements.
No end-to-end recognition gain on the actual-diagram controls is established.
The raw measurements are in \code{fast/results/graded\_transfer\_20261008.json}.


The conditional complexity target remains explicit.  If all allocated
intermediate state counts satisfy $N=n^{O(\log n)}$, all frontiers have
$k=O(\log^2n)$, and there are polynomially many stages, these explicit
operations fit a quasi-polynomial bound.  The integration proves neither
universal hypothesis.  Section~\ref{sec:rational-arithmetic} shows why even a
four-ended frontier can have large necessary expanded multiplicity.  The next
algebraic refinement in report 24 prunes paths irrelevant to survivors and
chooses propagation direction; its performance and topology obligations
remain separate from this kernel's correctness.
